%% Extraction schema as an appendix block. Paste at the end of Appendix A
%% (after Table tbl:vocab, before \section*{Acknowledgment}). Table tbl:schema
%% in Section 3.2 stays as the compact field summary; add this pointer to the
%% end of the "Schema" paragraph there:
%%   Appendix~\ref{app:schema} gives the record type as a listing together
%%   with one extracted chain from the corpus.
%%
%% Preamble additions:
%%   \usepackage{listings}
%%   \lstdefinestyle{schema}{basicstyle=\ttfamily\footnotesize,columns=fullflexible,
%%     keepspaces=true,breaklines=true,breakatwhitespace=true,frame=tb,
%%     framerule=0.4pt,xleftmargin=0pt,aboveskip=4pt,belowskip=4pt,
%%     comment=[l]{//},commentstyle=\color{gray}\itshape,showstringspaces=false}
%% Appendix figure numbering (the appendix currently resets tables only; add
%% these two lines next to the existing \setcounter{table}{0}):
%%   \setcounter{figure}{0}
%%   \renewcommand{\thefigure}{A.\arabic{figure}}
%%
%% Generated 2026-09-05 from event_extraction/out/full_corpus_v4.jsonl,
%% record 20020225X00252_1; the chain is reproduced verbatim.

Figure~\ref{fig:schema} states the record type behind Table~\ref{tbl:schema},
and Figure~\ref{fig:example} shows one chain as the extractor returned it.
The identifiers are the aggregation variables of the causation graph and
match the released output verbatim. Every field of the record
type is a closed vocabulary except the two free-text fields, the trigger
phrase and the evidence span, and both are constrained to the narrative. The
four rules that the grammar cannot express are enforced by the runner, which
requests one repair turn on violation; across the full corpus no record
required more than one. The JSON Schema
(\texttt{schema\_v4.json}), the system prompt (\texttt{system\_v4.txt}) and
the few-shot example (\texttt{few\_shot\_v4.json}) are released with the
code and match this listing verbatim.

\begin{figure}[H]
\centering
\begin{lstlisting}[style=schema]
Record {
  chain             : ChainNode[1..24]   // ordered: earliest factor first, outcome last
  outcome_severity  : fatal | serious_injury | minor_injury
                    | aircraft_damage_only | no_damage_or_injury | unknown
}
ChainNode {
  idx               : int                // equals the node's array position
  class             : condition | event | outcome
  factor_type       : one of 57 values   // 35 events, 12 conditions, 9 outcomes, unknown
  trigger           : string             // phrase from the narrative naming the factor
  phase_of_flight   : one of 22 values   // CICTT phase taxonomy, preflight .. post_impact
  cause_role        : primary | contributing | outcome | context
  caused_by         : Link[0..4]         // causal parents; empty for roots and conditions
}
Link {
  src               : int, src < idx     // backward only: every chain is acyclic
  strength          : direct | contributing   // mirrors the NTSB cause / factor flags
  evidence          : string             // narrative span supporting this link
}
// Enforced by the runner, outside the grammar: idx equals position; every src < idx;
// an outcome node is a parent only of outcome nodes; repeats of a factor_type merge.
\end{lstlisting}
\caption{The record type of the extraction schema. A record is one chain
of at most 24 nodes plus a record-level severity; causal structure lives in
each node's parent list. The 57 factor values are those of
Table~\ref{tbl:vocab}.}
\label{fig:schema}
\end{figure}

\begin{figure}[H]
\centering
\begin{lstlisting}[style=schema]
{"chain": [
 {"idx": 0, "class": "condition", "factor_type": "CARBURETOR_OR_INDUCTION_ICING",
  "trigger": "temperature 3C, dewpoint -6C, falling in range of susceptibility for carburetor icing",
  "phase_of_flight": "cruise", "cause_role": "contributing",
  "caused_by": []},
 {"idx": 1, "class": "event", "factor_type": "PROCEDURE_NOT_FOLLOWED",
  "trigger": "delayed use of carburetor heat",
  "phase_of_flight": "cruise", "cause_role": "primary",
  "caused_by": [
   {"src": 0, "strength": "contributing", "evidence": "carburetor icing conditions"}
  ]},
 {"idx": 2, "class": "event", "factor_type": "ENGINE_FAILURE",
  "trigger": "loss of engine power in cruise flight",
  "phase_of_flight": "cruise", "cause_role": "primary",
  "caused_by": [
   {"src": 1, "strength": "direct", "evidence": "delayed use of carburetor heat"}
  ]},
 {"idx": 3, "class": "condition", "factor_type": "UNSUITABLE_TERRAIN_FOR_FORCED_LANDING",
  "trigger": "snow covered field with powerlines and brush",
  "phase_of_flight": "emergency_descent", "cause_role": "contributing",
  "caused_by": []},
 {"idx": 4, "class": "outcome", "factor_type": "GROUND_IMPACT",
  "trigger": "struck a powerline and brush during the landing attempt",
  "phase_of_flight": "emergency_descent", "cause_role": "outcome",
  "caused_by": [
   {"src": 2, "strength": "direct", "evidence": "loss of engine power"},
   {"src": 3, "strength": "contributing", "evidence": "unsuitable terrain encountered"}
  ]},
 {"idx": 5, "class": "outcome", "factor_type": "INJURY_OR_FATALITY",
  "trigger": "pilot not injured",
  "phase_of_flight": "post_impact", "cause_role": "outcome",
  "caused_by": [
   {"src": 4, "strength": "direct", "evidence": "forced landing on road"}
  ]}
 ],
 "outcome_severity": "aircraft_damage_only"}
\end{lstlisting}
\caption{One extracted chain, NTSB record 20020225X00252: carburetor
icing conditions at the recorded temperature and dewpoint, delayed use of
carburetor heat, power loss in cruise, and a forced landing on unsuitable
terrain. Conditions enter with empty parent lists; every parent index is
smaller than the node's own; the two link strengths follow the NTSB cause and
factor flags; every link quotes the narrative span that supports it.}
\label{fig:example}
\end{figure}
